%!TEX root=../main.tex
\section{Introduction}
% Machine learning is being used increasingly to assess risk in various domains,
% including the criminal justice system, online advertising, and
% medical testing \cite{whitehouse-big-data}.
% %Fundamentally, a tool for assessing risk can be viewed as assigning
% %each individual an estimate of the probability that they belong
% %to the positive class in a classification problem.
% As with any application of machine learning, these risk assessments
% will be carried out with some underlying level of error.

Recently, there has been growing concern about errors of machine learning algorithms in sensitive domains -- including criminal justice, online advertising, and
medical testing \cite{whitehouse-big-data} -- which may systematically discriminate against
particular groups of people \cite{barocas2016big,berk2017fairness,chouldechova2017fair}.
%
%
% Example time
% \kilian{ I am not sure about this example, not everybody may agree.}
A recent high-profile example of these concerns was raised by the
news organization ProPublica, who studied a risk-assessment tool
that is widely used in the criminal justice system.
This tool assigns to each criminal defendant an estimated probability
that they will commit a future crime.
ProPublica found that the risk estimates assigned to
defendants who did not commit future crimes were on average
higher among African-American defendants than Caucasian defendants \cite{angwin-propublica-risk-scores}.
%Assigning non-trivial risk to people who didn't in fact commit a
%crime in the future
This is a form of false-positive error, and in this case it
disproportionately affected African-American defendants.
%
To mitigate issues such as these,
the machine learning community has proposed different
frameworks that attempt to quantify fairness in classification~\cite{barocas2016big,berk2017fairness,chouldechova2017fair,hardt2016equality,kleinberg2016inherent,woodworth2017learning,zafar2017fairness}.
A recent and particularly noteworthy framework is Equalized Odds \cite{hardt2016equality}
(also referred to as Disparate Mistreatment \cite{zafar2017fairness}),\footnote{
  For the remainder of the paper, we will use \emph{Equalized Odds}
  to refer to this notion of non-discrimination.
}
which constrains classification algorithms such that no error type (false-positive
or false-negative) disproportionately affects any population subgroup.
This notion of non-discrimination is feasible in many settings,
and researchers have developed tractable algorithms for achieving it \cite{hardt2016equality,goh2016satisfying,zafar2017fairness,woodworth2017learning}.
%
%These include trying to ensure \emph{individual fairness}
%\cite{dwork2012fairness}; training classifiers to
%guarantee fairness with respect to \emph{groups}, where a group may refer to a
%particular race, gender, or other subset of the population
%\cite{calders2009building,kamiran2009classifying,calders2010three}; and obfuscating the
%feature representation of individuals to remove any information of group
%membership \cite{zemel2013learning, louizos2015variational,
%edwards2015censoring}.

%\noindent\textbf{Calibration.}
%In this paper we investigate the relationship of non-discrimination and fairness.
%However, equalized error rates do not necessarily remove all potential for discrimination.
When risk tools are used in practice, a key goal is that they are
{\em calibrated}: if we look at the set of people who receive a predicted probability
of $p$, we would like a $p$ fraction of the members of this set to
be positive instances of the classification problem \cite{dawid1982well}.
Moreover, if we are concerned about fairness between two groups
$G_1$ and $G_2$ (e.g. African-American defendants and white defendants)
then we would like this calibration condition
to hold simultaneously for the set
of people within each of these groups as well
%: of all people in group $G_t$
%who are labeled $p$ (for $t = 1, 2$), a $p$ fraction should be
%positive instances
\cite{flores-re-propublica-fair}.
%
% Calibration is in a sense a necessary condition for a risk tool if it
% is going to be used as intended in practice.
Calibration is a crucial condition for risk tools in many settings.
If a risk tool for
evaluating defendants were not calibrated with respect to groups
defined by race, for example, then a probability estimate of $p$
could carry different meaning for African-American and white defendants,
and hence the tool would have the unintended and highly undesirable
consequence of incentivizing judges to take race into account when
interpreting its predictions.
% \footnote{There are many similar examples.
% Suppose a tool for estimating the probability that a medical student
% will succeed as a doctor were uncalibrated with respect to gender,
% and suppose this tool were used to inform hospital hiring decisions.
% Then it could become rational for patients choosing a doctor to take the
% doctor's gender into consideration as a criterion in their choice.}
%
Despite the importance of calibration as a property,
our understanding of how it interacts
with other fairness properties is limited.
We know from recent work that, except in the most constrained
cases, it is impossible to achieve calibration while also
satisfying Equalized Odds
\citep{kleinberg2016inherent,chouldechova2017fair}.
However, we do not know how best to achieve relaxations of these guarantees
that are feasible in practice.


%Our goal is to further investigate the relationship
% between calibration and equalized error rates.
% Our results generalize the work of \citet{kleinberg2016inherent}
% and \citet{chouldechova2017fair}, who first observed the incompatibility
% between calibration and Equalized Odds.
Our goal is to further investigate the relationship between calibration and error rates.
We show that even if the Equalized Odds conditions are relaxed
substantially -- requiring only that weighted sums of the group error
rates match -- it is still problematic to also enforce calibration.
We provide necessary and sufficient conditions under which this calibrated
relaxation
is feasible.
When feasible, it has a unique optimal solution that
can be achieved through post-processing of existing classifiers.
%
Moreover, we provide a simple post-processing algorithm to find this solution:
withholding predictive information for randomly chosen inputs
to achieve parity and preserve calibration. However, this simple
post-processing method is fundamentally unsatisfactory: although the
post-processed predictions of our information-withholding algorithm are ``fair'' in
expectation, most practitioners would object to the fact that a non-trivial
portion of the individual predictions are withheld as a result of coin tosses -- especially in
sensitive settings such as health care or criminal justice. The
optimality of this algorithm thus has troubling implications and shows that
calibration and error-rate fairness are inherently at odds (even beyond the initial results
by \cite{chouldechova2017fair} and \cite{kleinberg2016inherent}).

Finally, we evaluate these theoretical findings empirically,
comparing calibrated notions of non-discrimination against
the (uncalibrated) Equalized Odds framework on several datasets.
These experiments further support our conclusion that calibration and error-rate constraints are in most cases mutually incompatible goals.
In practical settings, it may be advisable to choose only one of these goals rather than attempting to achieve some relaxed notion of both.


% We characterize the set of calibrated classifiers for different population
% groups. With the help of a surprisingly simple geometric characterization we can show that

% . From this intuition, we can derive a number of straightforward
% yet profound implications.
% %
% 1. while there are notions of error-rate non-discrimination that are compatible
% with calibration, we show that such instances require a significant
% relaxation of the equalized odds conditions.
% 2. any such notions become infeasible when the prediction task
% is inherently noisy, or when population group membership is a strong predictive
% feature.
% 3. we demonstrate that achieving calibration and even the most simple
% non-discrimination constraints inherently incurs a non-trivial performance
% cost, which may be inadvisable in many situations.
% In the supplementary material, we show that all of these findings hold in
% an approximate sense as well; the (in)feasibility and performance
% costs degrade smoothly as the calibration condition is relaxed.
% We demonstrate these findings both theoretically and empirically,
% comparing calibrated notions of non-discrimination against
% the (uncalibrated) Equalized Odds framework on several datasets.
% Our conclusion is that calibration and error-rate constraints are in
% most cases mutually incompatible goals, and therefore practitioners
% should choose only one of these goals rather than attempting to achieve
% some relaxed notion of both.

% \gp{
% There is a unique solution
% can be achieved by a simple interpolation algorithm
% is unadvisable in many situations.
% }

% \gp{
% Can't do better than a random coin toss.
% }

% \gp{
% There has been a cry for fairness.
% }

% \gp{
% Here's a natural way of thinking about fairness.
% This simple geometric intuition makes it crystal clear
% that this notion is simply achievable, but unsatisfying because of coin toss.
% }

% \gp{
% Get rid of \methodname{}.
% }
